% !TEX root = ../main.tex
\section{From Model Output to Institutional Warrant}
\label{sec:warrant}

The account so far is epistemological, and it stops short of the question that
motivates it in practice: under what conditions may an institution justifiably
rely on a machine-generated claim? I close by stating the transition, not by
making it, since the constructive theory belongs to a companion paper.

The answer is not ``when the model knows.'' Sections~\ref{sec:output-operator}
and~\ref{sec:hyperintensionality} rule that out as a question, and
Section~\ref{sec:internal} explains why even a truth-correlated internal
representation would not supply it. A raw output is better understood as
\emph{unlicensed testimony}: a candidate proposition requiring evaluation. That a
model produced $\phi$ does not authorise reliance on $\phi$.

What can bear the weight is warrant, and warrant is a property of a governed
process rather than of model parameters. A preliminary schema:
\begin{equation}
  \Warranted_{I,c}(\phi)
  \;\Leftrightarrow\;
  \Resolved(\phi, c) \wedge \Grd_M(\phi \mid c)
  \wedge \Verified(V, \phi, c) \wedge \Authorized(I, \phi, c),
  \label{eq:warrant}
\end{equation}
for an institution $I$, context $c$, and verification procedure $V$. The schema is
schematic, as the notation of this paper generally is: it states necessary
components rather than a reduction, and each conjunct names a body of work rather
than a settled condition.

Three things are worth noting about how the present account feeds it.

First, $\Resolved(\phi, c)$ is not an administrative preliminary. By
Section~\ref{sec:contextual-truth}, an unresolved claim is not yet evaluable, so
resolution is prior to grounding rather than parallel to it. In deployment this
is the step most often skipped, because natural language does not announce that
its decisive parameters are missing.

Second, $\Grd_M(\phi \mid c)$ is doing the work that ``justification'' does in the
doxastic case, and Section~\ref{sec:measurement} showed it is measurable only for
architectures that make $E$ explicit and applicable at the point of assertion.
Warrant therefore has an architectural precondition, not merely a procedural one.

Third, verification and authorization are irreducible to either. A grounded,
resolved, correctly asserted claim may still be one an institution is not
entitled to act on, and no amount of improvement to $M$ changes that. This is why
the sequence
\[
  \text{raw output} \rightarrow \text{candidate claim} \rightarrow
  \text{grounded claim} \rightarrow \text{verified claim} \rightarrow
  \text{authorized reliance}
\]
is a sequence of distinct conditions rather than a gradient of confidence. The
model may contribute at every stage---retrieving, summarising, comparing,
structuring, explaining. Its output is not identical with the warranted
conclusion at any of them.
